% ============================================================================
% Fichier  : front/01_avant_propos.tex
% Titre    : Avant-propos
% Auteur   : MINKA MI NGUIDJOÏ Thierry Emmanuel
% Version  : 1.0 -- novembre 2026
% ============================================================================
% Formule de citation de la méthode : cadre de cohérence du corpus (A0 pour
% le cadre, A1 pour l'instrument ; A2 à A5 « en cours de formalisation »).
% Le nom d'usage « Minka-N8 » n'est pas employé (non tranché par l'auteur).
% ============================================================================

\chapter*{Avant-propos}
\addcontentsline{toc}{chapter}{Avant-propos}
\markboth{Avant-propos}{Avant-propos}

\epigraph{« Ce que l'organisation ignore de ses propres systèmes conditionne
la qualité des décisions qu'elle peut prendre à leur sujet. »}%
{\textit{Vers l'auditologie}~\cite{minka2026auditologie}}

Ce cours s'adresse aux étudiants du Master Professionnel en Management des
Systèmes Numériques de la Chaire UNESCO d'Afrique Centrale Accès TIC. Il
prolonge l'UE MSN-2S3-MSI01 (Gouvernance stratégique des SI) et dialogue
avec l'UE MSN-2S3-MSI03 (IT Governance, COBIT et ITIL).

Il repose sur une double conviction. La première est qu'auditer ne consiste
pas à remplir une grille, mais à réduire l'incertitude d'un tiers sur un
système qu'il ne peut pas observer seul. Un conseil d'administration, un
régulateur, un juge, un commissaire aux comptes s'appuient sur ce que
l'auditeur affirme~; ils doivent pouvoir savoir aussi ce qu'il ignore.
C'est pourquoi ce cours admet, à côté de \OUI{} et de \NON{}, une troisième
réponse~: \NEK{}, « connaissance insuffisante ».

La seconde conviction est que le métier ne s'importe pas tel quel. Les
normes de l'IIA, de l'ISACA, de l'ISO et de l'INTOSAI sont des acquis
précieux~; elles s'appliquent ici sous le regard de la loi camerounaise sur
la cybersécurité, de la COBAC, du droit OHADA et de la Chambre des comptes,
dans des organisations dont les moyens et les contraintes sont les nôtres.

\bigskip

\textbf{Ce que ce cours n'est pas.} Il n'est pas un résumé des normes
professionnelles, ni une préparation à une certification, ni un cours de
tests d'intrusion. Les normes y sont citées avec exactitude, dans l'édition
en vigueur quand elle a pu être vérifiée~; les textes de loi n'y sont
cités entre guillemets que s'ils ont été lus dans leur version officielle.

\bigskip

\textbf{Ce que ce cours est.} Un traité en six parties, précédées d'un
prologue, qui va du cas concret vers la méthode, de la méthode vers la
pratique, et de la pratique vers la recherche. La Partie~I pose les
fondements~: audit, contrôle interne, risque, fraude. La Partie~II donne la
méthode de la mission, de la lettre de mission au rapport. La Partie~III
présente les normes, les référentiels et le droit. La Partie~IV audite le
SI domaine par domaine, et se termine par trois cas intégrateurs. La
Partie~V aborde l'audit continu et l'intelligence artificielle. La
Partie~VI présente le programme de recherche sur lequel la méthode
s'appuie.

\bigskip

\textbf{Un cours, un corpus.} La méthode enseignée est l'instanciation
opérationnelle, pour le praticien, du programme de recherche de
l'auditologie~\cite{minka2026auditologie}~; son instrument de réponse
\OUI{}~/ \NON{}~/ \NEK{} et ses mesures viennent de l'article
\cite{minka2026nek}. Ces deux textes sont des prépublications déposées sur
HAL. La méthode est présentée intégralement dans ce cours~; la
formalisation de la décomposition, du suivi longitudinal, du signalement
stratégique et de l'instrumentation adaptative est en cours de
publication, et le cours la signale comme telle.

\bigskip

\textbf{Trois organisations fictives.} Le cours suit trois missions.
\emph{LiZbiZ}, entreprise de logistique aux quatre centres de données,
reprise de l'article fondateur, est le cas principal~: le cabinet AUDIX
CONSEIL y conduit deux cycles d'audit. La \emph{BANQUE SANAGA}, à Douala,
affronte les suites d'une panne et le regard de la COBAC. L'\emph{ANPP},
agence publique de paiement, est auditée par la Chambre des comptes. Les
noms, les dates et les montants sont fictifs~; les situations sont
plausibles, et chaque chiffre se recalcule.

\bigskip

\textbf{Comment travailler.} Lisez chaque chapitre avec sa fiche
d'auto-évaluation~; refaites les calculs des lectures auditologiques~;
faites les labs de NDAP SUKULU avant de lire leurs corrigés en séance.
Surtout, prenez l'habitude que le cours cherche à transmettre~: séparer ce
que l'on a constaté de ce que l'on en pense, et dire ce que l'on ne sait
pas.

\bigskip

\begin{flushright}
    \textit{\auteurcours}\\
    \textit{Yaoundé, novembre 2026}
\end{flushright}
